\subsection{Planitud y dimensión.}

\begin{proposition}[6.1.1]
\label{IV.2.s6.1.1-fr}
Sean $A$, $B$ dos anillos locales noetherianos, $\mathfrak m$ el ideal maximal
de $A$, $k=A/\mathfrak m$ su cuerpo residual,
$\varphi:A\to B$ un homomorfismo local. Se supone que, para todo ideal
primo $\mathfrak p$ de $A$, distinto de $\mathfrak m$, y para todo elemento
minimal $\mathfrak q$ del conjunto de los ideales primos de $B$ que contienen
$\mathfrak pB$, $\varphi^{-1}(\mathfrak q)$ sea distinto de $\mathfrak m$
\textup{(dicho de otro modo, que ninguna componente irreducible de
$\operatorname{Spec}(B/\mathfrak pB)$ esté contenida en la imagen inversa
del punto cerrado $\mathfrak m$ de $\operatorname{Spec}(A)$ en
$\operatorname{Spec}(B)$)}. Entonces se tiene
\begin{equation}
\dim(B)=\dim(A)+\dim(B\otimes_A k).
\tag{6.1.1.1}\label{IV.2.s6.1.1.1-fr}
\end{equation}
\end{proposition}

Razonemos por inducción sobre $n=\dim(A)$; la afirmación es trivial para
$n=0$, pues $\mathfrak mB$ está entonces contenido en el nilradical de $B$, ya
que $\mathfrak m$ es el nilradical de $A$. Se puede, pues, suponer $n>0$. Sean
$\mathfrak q_i$ $(1\leq i\leq m)$ los ideales primos minimales de $B$, y sea
$\mathfrak p_i=\varphi^{-1}(\mathfrak q_i)$; para todo $i$, se tiene
$\mathfrak p_i\neq\mathfrak m$; de lo contrario existiría (puesto que $n>0$)
un ideal primo $\mathfrak p\neq\mathfrak m$ contenido en $\mathfrak p_i$ y,
como $\mathfrak q_i$ es minimal entre los ideales primos de $B$ que contienen
$\mathfrak pB$, se llegaría a una contradicción con la hipótesis. Por tanto,
$\mathfrak m$ es distinto de la unión de los $\mathfrak p_i$ y de los ideales
minimales $\mathfrak p'_j$ $(1\leq j\leq r)$ de $A$ (Bourbaki,
\emph{Alg. comm.}, cap.~II, §~1, n\textsuperscript{o}~1, prop.~2), y existe
un $a\in\mathfrak m$ que no pertenece a ninguno de los $\mathfrak p_i$ ni de
los $\mathfrak p'_j$. Pongamos $A'=A/aA$, $B'=B/aB$; se tiene (\textbf 0, 16.3.4)
\[
\dim(A')=\dim(A)\setmin 1,\qquad \dim(B')=\dim(B)\setmin 1
\]
por construcción de $a$, y, por otra parte,
$B'\otimes_{A'}k=B\otimes_Ak=B/\mathfrak mB$, por tanto
\[
\dim(B'\otimes_{A'}k)=\dim(B\otimes_Ak)
\]
y basta, en consecuencia, demostrar (6.1.1.1) para $A'$ y $B'$; ahora bien, en
virtud de la correspondencia biunívoca entre los ideales de $A$ (resp. $B$)
que contienen $a$ y los ideales de $A'$ (resp. $B'$), las hipótesis del
enunciado son también válidas para $A'$ y $B'$; se puede, pues, aplicar la
hipótesis de inducción, lo que termina la demostración.

\begin{corollary}[6.1.2]
\label{IV.2.s6.1.2-fr}
Sean $A$, $B$ dos anillos locales noetherianos, $k$ el cuerpo residual de
$A$, $\varphi:A\to B$ un homomorfismo local, $M\neq0$ un $A$-módulo de tipo
finito, y $N\neq0$ un $B$-módulo de tipo finito. Si $N$ es un $A$-módulo plano
\textup{(resp.} si $B$ es un $A$-módulo plano), se tiene
\begin{equation}
\dim_B(M\otimes_A N)=\dim_A(M)+
\dim_{B\otimes_Ak}(N\otimes_Ak)
\tag{6.1.2.1}\label{IV.2.s6.1.2.1-fr}
\end{equation}
\textup{(resp. (6.1.1.1))}.
\end{corollary}

Basta demostrar la afirmación relativa a $N$; por otra parte, si
$\mathfrak b$ es el anulador de $N$, se puede sustituir $B$ por
$B/\mathfrak b$ y, por tanto, suponer que
$\operatorname{Supp}(N)=\operatorname{Spec}(B)$; la hipótesis significa
entonces que el morfismo $f:\operatorname{Spec}(B)\to\operatorname{Spec}(A)$
correspondiente a $\varphi$ es \emph{cuasiplano} (2.3.3); se concluye
(2.3.4) que, para todo ideal primo $\mathfrak p\neq\mathfrak m$ de $A$,

\oldpage[IV-2]{136}

toda componente irreducible de $f^{-1}(V(\mathfrak p))$ \emph{domina}
$V(\mathfrak p)$ y, por consiguiente, se satisface la condición de (6.1.1),
de donde resulta la conclusión.
